\handout
{Matrix and System of Equations}
{\textsc{Linear Algebra (LA)}}
{\href{https://creativecommons.org/publicdomain/zero/1.0/}{CC0 1.0 Universal license}}
{Author: Jaehoon Song}{Release: \generatedOn}
{LA: Matrix and System of Equations}



\begin{enumerate}
  \item \textbf{Matrix Equation by Inverse Matrix} (9 points)\\
  Use the inverse matrix method to solve the following systems of equations.
  \begin{enumerate}
    \item Solve for the matrix \(X\) in the equation shown. (3 points)
    \[
    X\begin{bmatrix}
    -1 & 4 \\
    2 & 9
    \end{bmatrix} =
    \begin{bmatrix}
    -13 & -50 \\
    -8 & -19
    \end{bmatrix}.
    \]
    \begin{subanswer}
      your answer here.
    \end{subanswer}

    \item Solve for the matrix \(X\) in the equation shown. (3 points)
    \[
    X \begin{bmatrix} 6 & 1 \\ -2 & 4 \end{bmatrix} = \begin{bmatrix} -2 & 0 \\ 14 & -2 \end{bmatrix}
    \]
    \begin{subanswer}
      your answer here.
    \end{subanswer}

    \item Solve for the matrix \(X\) in the equation shown. (3 points)
    \[
    X\begin{bmatrix}
    1 & -2 \\
    4 & -9
    \end{bmatrix} =
    \begin{bmatrix}
    -5 & 12 \\
    9 & -20 \\
    4 & -8
    \end{bmatrix}.
    \]
    \begin{subanswer}
      your answer here.
    \end{subanswer}
  \end{enumerate}

  \newpage

  \item \textbf{Matrix Equation by RREF} (9 points)\\
  Solve the following systems of equations by writing a matrix to 
  represent each system and reducing it to \textbf{reduced row echelon form (RREF)}. 
  Include the row operations you use for decomposition.
  \begin{enumerate}
    \item Solve the system of equations. (3 points)
    \[
    \begin{aligned}
    2x + 4k - z &= 15, \\
    2z + k &= -2, \\
    4z + 3y - 3k &= -9, \\
    2z - y &= -11.
    \end{aligned}
    \]
    \begin{subanswer}
      your answer here.
    \end{subanswer}

    \item Solve the system of equations. (3 points)
    \[
    \begin{aligned}
    2x + y &= 2, \\
    y + z + 4k &= 21, \\
    2k - z &= 13, \\
    2z - y &= -10.
    \end{aligned}
    \]
    \begin{subanswer}
      your answer here.
    \end{subanswer}

    \item Solve the system of equations. (3 points)
    \[
    \begin{aligned}
    2x - 8z + k &= 14, \\
    -4z + 2y - 2 &= 0, \\
    4y - 9k &= -12, \\
    x + k - z &= 1.
    \end{aligned}
    \]
    \begin{subanswer}
      your answer here.
    \end{subanswer}
  \end{enumerate}














  \newpage

  \item \textbf{Matrix Equation from System of Equations} (4 points)\\
  Minju, Jiya, Hanna, and Jackie are ordering 
  candy and flower arrangements for their mothers for 
  Mother's Day. Each box of candy has a fixed cost, 
  and each flower arrangement has a fixed cost, but the 
  website they are ordering from does not advertise the price. 
  Jackie's mom is allergic to flowers and Hanna's mom does not 
  like candy, so they only want to order something their mom 
  likes. To figure out the cost, Minju and Jiya order first. 
  Minju orders 3 boxes of candy and 4 flower arrangements, and 
  her total is \(\$38\). Jiya orders 12 boxes of candy and 2 
  flower arrangements, and her total is \(\$40\). Hanna and 
  Jackie each have \(\$45\) to spend. Use this information 
  to answer the following questions:
  \begin{enumerate}
    \item What is the cost per unit for flower arrangements? (1 point)
    \begin{subanswer}
      your answer here.
    \end{subanswer}
    \item What is the cost per unit for boxes of candy? (1 point)
    \begin{subanswer}
      your answer here.
    \end{subanswer}
    \item How many boxes of candy can Jackie buy? (1 point)
    \begin{subanswer}
      your answer here.
    \end{subanswer}
    \item How many flower arrangements can Hanna buy? (1 point)
    \begin{subanswer}
      your answer here.
    \end{subanswer}
  \end{enumerate}
  












  \newpage
  \item \textbf{Matrix Transformation (Points)} (11 points)\\
  Answer the following questions about transformations of triangles:
  \begin{enumerate}
    \item A triangle has vertices \(A(1, 2)\), \(B(3, 4)\), and \(C(5, 6)\).
    The triangle is then rotated \(90^{\circ}\) counterclockwise around the origin.
    \begin{enumerate}
      \item Find the coordinates of \(A'\), the image of vertex \(A\) after the rotation. (1 point)
      \begin{subanswer}
        your answer here.
      \end{subanswer}
      \item Find the coordinates of \(B'\), the image of vertex \(B\) after the rotation. (1 point)
      \begin{subanswer}
        your answer here.
      \end{subanswer}
      \item Find the coordinates of \(C'\), the image of vertex \(C\) after the rotation. (1 point)
      \begin{subanswer}
        your answer here.
      \end{subanswer}
      \item Find the area of triangle \(A'B'C'\). (1 point)
      \begin{subanswer}
        your answer here.
      \end{subanswer}
    \end{enumerate}

    \item $\triangle ABC$ has vertices $A(a, 2)$, $B(-1, 3)$, and $C(6, a)$. 
    It is known that $\triangle ABC$ has an area of 6 square units.
    \begin{enumerate}
      \item Find all possible values of $a$. Show the work that leads to your answer. (3 points)
      \begin{subanswer}
        your answer here.
      \end{subanswer}
      \item If $\triangle ABC$ is rotated $15^\circ$ clockwise around the origin, find 
      the coordinates of $B'$, the image of vertex $B$ after the rotation. (4 points)\\
      (Hint: $\cos(a \pm b) = \cos(a)\cos(b) \mp \sin(a)\sin(b)$ and $\sin(a \pm b) = 
      \sin(a)\cos(b) \pm \cos(a)\sin(b)$)
      \begin{subanswer}
        your answer here.
      \end{subanswer}
    \end{enumerate}
  \end{enumerate}







  \newpage

  \item \textbf{Matrix Transformation (Equations)} (12 points)\\
  Various conic sections are given in parametric form. Each conic section is then 
  rotated $45^{\circ}$ counterclockwise around the origin. Use this information to 
  answer the following questions:
  \begin{enumerate}
    \item A parabola is given as $x(t)=-t \sqrt{3}, y(t)=-12 t^2$.
    \begin{enumerate}
      \item Find the equation of the \textbf{parabola} after the transformation. (2 points)
      \begin{subanswer}
        your answer here.
      \end{subanswer}
      \item Find the \textbf{focus} of the transformed parabola. (1 point)
      \begin{subanswer}
        your answer here.
      \end{subanswer}
      \item Find the equation of the \textbf{directrix} of the transformed parabola. (1 point)
      \begin{subanswer}
        your answer here.
      \end{subanswer}
    \end{enumerate}

    \item A circle is given as $x(t)=3\cos(t), y(t)=3\sin(t)$.
    \begin{enumerate}
      \item Find the equation of the \textbf{circle} after the transformation. (2 points)
      \begin{subanswer}
        your answer here.
      \end{subanswer}
      \item Find the \textbf{center} of the transformed circle. (1 point)
      \begin{subanswer}
        your answer here.
      \end{subanswer}
      \item Find the \textbf{radius} of the transformed circle. (1 point)
      \begin{subanswer}
        your answer here.
      \end{subanswer}
    \end{enumerate}

    \item An ellipse is given as $x(t)=5\cos(t), y(t)=3\sin(t)$.
    \begin{enumerate}
      \item Find the equation of the \textbf{ellipse} after the transformation. (2 points)
      \begin{subanswer}
        your answer here.
      \end{subanswer}
      \item Find the \textbf{foci} of the transformed ellipse. (1 point)
      \begin{subanswer}
        your answer here.
      \end{subanswer}
      \item Find the \textbf{major and minor axes} of the transformed ellipse. (1 point)
      \begin{subanswer}
        your answer here.
      \end{subanswer}
    \end{enumerate}

    \item A hyperbola is given as $x(t)=4\sec(t), y(t)=2\tan(t)$.
    \begin{enumerate}
      \item Find the equation of the \textbf{hyperbola} after the transformation. (2 points)
      \begin{subanswer}
        your answer here.
      \end{subanswer}
      \item Find the \textbf{foci} of the transformed hyperbola. (1 point)
      \begin{subanswer}
        your answer here.
      \end{subanswer}
      \item Find the equations of the \textbf{asymptotes} of the transformed hyperbola. (1 point)
      \begin{subanswer}
        your answer here.
      \end{subanswer}
    \end{enumerate}
  \end{enumerate}





\end{enumerate}
